\documentclass[10pt]{scrartcl}
\usepackage[mdthm,colorsec,noauthor,noasy]{evan}
\usepackage{handout}
\DeclareMathOperator{\im}{range}
\title{MATH 110 --- Week 4}
\subtitle{Notes Chapters 7--8; Homework 4}
\author{}\date{}\ihead{}
\begin{document}\maketitle
\section{Isomorphisms and inverses}\begin{qbox}When can a linear map be reversed? The map $T(x,y)=(x+y,y)$ has inverse $S(a,b)=(a-b,b)$, since both compositions return the input.\end{qbox}
\begin{prob}\begin{definition*}[Notes Definition 21; LADR 3.59, 3.69]A linear map $T:V\to W$ is invertible, or an isomorphism, if a linear map $S:W\to V$ satisfies $ST=I_V$ and $TS=I_W$. Such an $S$ is an inverse. The spaces are isomorphic, written $V\cong W$, if an isomorphism exists. A left inverse satisfies $ST=I_V$; a right inverse satisfies $TS=I_W$.\end{definition*}\end{prob}
\begin{proposition*}[Notes Proposition 28; LADR 3.60]An inverse is unique, denoted $T^{-1}$. More strongly, any left inverse $L$ and right inverse $R$ of the same map satisfy $L=R$.\end{proposition*}
\begin{motive}\ii[Want.] Compare two candidate inverses without solving their formulas.\ii[Know.] Composition is associative, and the identities act trivially.\ii[Key step.] Insert $TR=I_W$ between $L$ and its input.\end{motive}
\begin{prob}\begin{proof}\begin{steps}\ii $L=LI_W=L(TR)=(LT)R=I_VR=R$.\ii Any two two-sided inverses are in particular a left and right inverse, so they agree.\end{steps}\end{proof}\end{prob}
\begin{example*}[Notes Chapter 7: one-sided nonuniqueness]Let $A:F^2\to F$ have row matrix $(1\ \ 2)$.
\begin{steps}\ii For every $t\in F$, the column $R_t=(1-2t,t)^{\mathsf T}$ gives $AR_t=1$. Different $t$ give different right inverses.
\ii There is no left inverse: $A(-2,1)=0$ for a nonzero vector, so $LA=I$ would imply $(-2,1)=L0=0$.
\ii There are infinitely many right inverses when $F$ is infinite, but only $|F|$ of this displayed family when $F$ is finite.\end{steps}\end{example*}
\begin{proposition*}[Notes Proposition 29; LADR 3.63]A linear map is invertible iff it is both injective and surjective. The inverse of a bijective linear map is automatically linear.\end{proposition*}
\begin{motive}\ii[Want.] Turn a set-theoretic inverse into a linear inverse.\ii[Know.] Bijectivity gives a unique preimage of every output.\ii[Key step.] Apply $T$ to the proposed linearity equation and use injectivity.\end{motive}
\begin{prob}\begin{proof}\begin{steps}\ii If $T^{-1}$ exists and $Tv=0$, then $v=T^{-1}Tv=T^{-1}0=0$, so $T$ is injective. Every $w=T(T^{-1}w)$ is in its range, so it is onto.\ii Conversely, define $S(w)$ as the unique $v$ with $Tv=w$. Then $TS=I_W$ and $ST=I_V$.\ii For $a\in F$, $T(S(aw_1+w_2))=aw_1+w_2=T(aSw_1+Sw_2)$. Injectivity yields $S(aw_1+w_2)=aSw_1+Sw_2$, proving linearity.\end{steps}\end{proof}\end{prob}
\begin{proposition*}[Notes Proposition 30; LADR 3.70]Two finite-dimensional vector spaces over the same field are isomorphic iff their dimensions agree.\end{proposition*}
\begin{motive}\ii[Want.] Decide existence without a preassigned map.\ii[Know.] Rank--nullity and the unique map with prescribed basis images (Week 3).\ii[Key step.] Pair corresponding vectors in two equally long bases.\end{motive}
\begin{prob}\begin{proof}\begin{steps}\ii If $T:V\to W$ is an isomorphism, its kernel is zero and its range is $W$. Rank--nullity gives $\dim V=\dim W$.\ii If both dimensions are $n$, choose bases $v_1,\ldots,v_n$ and $w_1,\ldots,w_n$. Define linear maps by $Tv_j=w_j$ and $Sw_j=v_j$.\ii The compositions agree with the identity on their bases and therefore everywhere. Hence $T$ is an isomorphism.\end{steps}\end{proof}\end{prob}
\begin{proposition*}[Notes Chapter 7: infinite-dimensional variant]
Two spaces with bases indexed by the same set are isomorphic. Conversely, an isomorphism sends any basis to a basis.
\end{proposition*}
\begin{motive}\ii[Want.] Extend the basis-pairing construction to arbitrary bases.\ii[Know.] Every vector is a finite linear combination of basis vectors.\ii[Key step.] The construction needs only finite sums, regardless of the size of the full basis.\end{motive}
\begin{prob}\begin{proof}\begin{steps}
\ii Pair the basis vectors by their common index set and define $T$ on finite linear combinations using the same coefficients. Unique basis expansion makes it well-defined and linear; reversing the pairing gives an inverse.
\ii Conversely, if $T$ is an isomorphism, its images of a basis span by surjectivity. A finite relation among these images becomes a relation in the original basis after applying $T^{-1}$, so the images are independent.
\ii The existence of arbitrary bases and their cardinal comparison use the choice results cited in Week 2. The notation $\dim V=\infty$ alone does not specify the cardinality of a basis.
\end{steps}\end{proof}\end{prob}
\begin{proposition*}[Notes Proposition 31; LADR 3.65]For $\dim V=\dim W<\infty$ and linear $T:V\to W$, invertibility, injectivity, and surjectivity are equivalent.\end{proposition*}
\begin{motive}\ii[Want.] Prove invertibility using just one condition.\ii[Know.] Week 3 Proposition 25 makes injectivity and surjectivity equivalent in equal finite dimensions.\ii[Key step.] Combine that criterion with Proposition 29.\end{motive}
\begin{prob}\begin{proof}\begin{steps}\ii Either injectivity or surjectivity gives both by Week 3, hence invertibility by Proposition 29.\ii Invertibility gives both properties by Proposition 29, proving all implications.\end{steps}\end{proof}\end{prob}
\method{Compute an inverse by solving for the input}\begin{steps}\ii If $(a,b)=T(x,y)=(x+y,y)$, solve $y=b$ and $x=a-b$.\ii Thus $T^{-1}(a,b)=(a-b,b)$; substitution verifies $T^{-1}T=I$ and $TT^{-1}=I$.\end{steps}
\begin{shortcuts}\ii For equal finite dimensions, prove zero kernel or full range; Proposition 31 supplies invertibility.\ii To prove two linear maps equal, compare them only on a basis (Week 3).\end{shortcuts}\begin{trapbox}\ii A one-sided inverse need not be unique or two-sided for unequal dimensions: the row $(1\ \ 2)$ above has multiple right inverses and no left inverse.\ii Two abstractly isomorphic spaces do not make every map between them invertible: the zero map on $\mathbb R$ is a counterexample.\end{trapbox}
\section{The two-by-two formula and polynomial existence}\begin{qbox}Can we decide invertibility before solving every output equation? For a two-by-two matrix, one scalar controls an explicit inverse formula.\end{qbox}
\begin{prob}\begin{definition*}[LADR 3.79--3.80]A square matrix $A\in F^{n,n}$ is invertible if there is $B\in F^{n,n}$ with $AB=BA=I_n$; its unique inverse is denoted $A^{-1}$. This is exactly invertibility of the map $x\mapsto Ax$, by Week 3 matrix composition. \end{definition*}\end{prob}
\begin{prob}\begin{definition*}[Notes Chapter 7: two-by-two determinant]For $A=\left(\begin{smallmatrix}a&b\\c&d\end{smallmatrix}\right)$ define $\det A=ad-bc$. This notation here is only the two-by-two formula; no general determinant theory is assumed.\end{definition*}\end{prob}
\begin{proposition*}[Notes Chapter 7; HW4 P1]$A$ is invertible iff $ad-bc\ne0$. In that case $\displaystyle A^{-1}=\frac1{ad-bc}\begin{pmatrix}d&-b\\-c&a\end{pmatrix}$.\end{proposition*}
\begin{motive}\ii[Want.] Produce a formula and prove the zero case is impossible.\ii[Know.] Entries lie in a field, so a nonzero scalar can be divided out; invertible means bijective.\ii[Key step.] Multiply by the matrix with swapped diagonal entries and negated off-diagonal entries.\end{motive}
\begin{prob}\begin{proof}\begin{steps}\ii For $B=\left(\begin{smallmatrix}d&-b\\-c&a\end{smallmatrix}\right)$, direct multiplication gives $AB=BA=(ad-bc)I_2$: the off-diagonal entries cancel, and each diagonal entry is $ad-bc$.\ii If $ad-bc\ne0$, scaling $B$ by its reciprocal gives both inverse identities.\ii If $ad-bc=0$ and $(a,b)\ne(0,0)$, the nonzero vector $(b,-a)$ satisfies $A(b,-a)=(0,cb-da)=0$. Thus $A$ is not injective.\ii If $a=b=0$, the first coordinate of every output is zero, so $A$ is not onto. Proposition 29 excludes invertibility in either case.\end{steps}\end{proof}\end{prob}
\begin{proposition*}[Polynomial root bound used in Notes Example 34]A nonzero real polynomial of degree at most $n$ has at most $n$ distinct real roots.\end{proposition*}
\begin{motive}\ii[Want.] Use interpolation conditions to force a polynomial to vanish.\ii[Know.] For $k\ge1$, $x^k-a^k=(x-a)\sum_{j=0}^{k-1}x^{k-1-j}a^j$.\ii[Key step.] Factor out each root using this identity; different remaining roots survive in the quotient.\end{motive}
\begin{prob}\begin{proof}\begin{steps}\ii If $p(a)=0$, write $p(x)=p(x)-p(a)$ and factor each monomial difference using the displayed identity. Thus $p=(x-a)q$, with degree reduced by one for nonzero $p$.\ii If $b\ne a$ is another root, $0=p(b)=(b-a)q(b)$ implies $q(b)=0$. Repeat for distinct roots.\ii After $r$ distinct roots have been factored, the degree is at least $r$. Hence $r\le n$.\end{steps}\end{proof}\end{prob}
\begin{example*}[Notes Example 34: interpolation]Let $x_0,\ldots,x_n$ be distinct real numbers and $y_0,\ldots,y_n$ arbitrary.
\begin{steps}\ii Define $T:\mathcal P_n(\mathbb R)\to\mathbb R^{n+1}$ by $Tp=(p(x_0),\ldots,p(x_n))$. Evaluation is linear.
\ii If $Tp=0$, the root bound forces $p=0$. Both spaces have dimension $n+1$, so Proposition 31 gives bijectivity.
\ii Thus a unique polynomial satisfies all $p(x_j)=y_j$. A computing formula is $p(x)=\sum_{j=0}^n y_j\prod_{k\ne j}(x-x_k)/(x_j-x_k)$: at $x_i$, every summand is zero except $j=i$, which equals $y_i$.
\ii For points $(0,1),(1,2),(2,5)$, the formula gives $\frac12(x-1)(x-2)-2x(x-2)+\frac52x(x-1)=1+x^2$.\end{steps}\end{example*}
\begin{example*}[Notes Example 35; LADR 3.67]Given any $q\in\mathcal P(\mathbb R)$, there exists a polynomial $p$ with $((x^2+5x+7)p)''=q$.
\begin{steps}\ii Choose $m\ge0$ with $q\in\mathcal P_m(\mathbb R)$ and define $T(p)=((x^2+5x+7)p)''$ on $\mathcal P_m(\mathbb R)$. This is linear; a nonzero degree-$r$ input produces degree $r$ output by the power rule, so the map stays in this space.
\ii If $Tp=0$, the product $(x^2+5x+7)p$ has second derivative zero, hence degree at most $1$ or is zero. A nonzero $p$ would give degree at least $2$, impossible. Thus $p=0$.
\ii Proposition 31 gives surjectivity, proving existence. For example $q=2$ has solution $p=1$ by differentiating $x^2+5x+7$ twice.\end{steps}\end{example*}
\begin{example*}[Notes Example 36]On the infinite-dimensional space $\mathcal P(\mathbb R)$, differentiation is onto but not injective.
\begin{steps}\ii Every $q=\sum_{j=0}^m a_jx^j$ is the derivative of $\sum_{j=0}^m a_jx^{j+1}/(j+1)$.
\ii Nonzero constants have derivative zero, so the kernel is nonzero. Thus the equal finite-dimension criterion cannot be applied here.\end{steps}\end{example*}
\begin{shortcuts}\ii Two-by-two invertibility needs only $ad-bc$; over a finite field reduce all arithmetic in that field.\ii An equal-dimensional polynomial map may be easier to prove injective by degree than to solve coefficient by coefficient.\ii For interpolation, use the product formula to compute, and the root bound to prove uniqueness.\end{shortcuts}\begin{trapbox}\ii Distinct interpolation inputs are necessary for arbitrary outputs: $p(0)=1$ and $p(0)=2$ cannot both hold.\ii The two-by-two inverse requires division by a nonzero determinant: $\left(\begin{smallmatrix}1&2\\2&4\end{smallmatrix}\right)$ has nonzero entries but is not invertible.\end{trapbox}
\section{Products and one-sided inverses}\begin{qbox}Why is the inverse of a product written in reversed order? To reverse $ST$, first reverse the final action $S$, then the action $T$.\end{qbox}
\begin{proposition*}[Notes Proposition 32; LADR 3D.2, 3D.11]If operators $S,T$ are invertible, $(ST)^{-1}=T^{-1}S^{-1}$. If $ST$ is invertible, then $S$ is surjective and $T$ is injective. In finite dimension the latter condition also makes both $S,T$, and therefore $TS$, invertible.\end{proposition*}
\begin{motive}\ii[Want.] Extract information about factors from their product.\ii[Know.] The inverse identities and associativity are available.\ii[Key step.] For necessity, use an arbitrary target for $S$ and a kernel vector for $T$.\end{motive}
\begin{prob}\begin{proof}\begin{steps}\ii $(T^{-1}S^{-1})(ST)=T^{-1}(S^{-1}S)T=I$ and $(ST)(T^{-1}S^{-1})=S(TT^{-1})S^{-1}=I$. This proves the inverse formula.\ii If $ST$ is invertible, any $v=ST((ST)^{-1}v)$ is an image under $S$, so $S$ is onto.\ii If $Tv=0$, then $STv=0$; injectivity of $ST$ gives $v=0$. Thus $T$ is injective.\ii In finite dimension Proposition 31 makes both factors invertible, and the first part also applies to $TS$.\end{steps}\end{proof}\end{prob}
\begin{proposition*}[Notes Chapter 7: restricted factors]Even without finite dimension, if $ST$ is invertible, then $T:V\to\im T$ and $S|_{\im T}:\im T\to V$ are isomorphisms, and $(ST)^{-1}=T^{-1}(S|_{\im T})^{-1}$, with these restricted domains understood.\end{proposition*}
\begin{motive}\ii[Want.] Invert the factors on the spaces where they are bijective.\ii[Know.] Proposition 32 gives injectivity of $T$ and surjectivity of $ST$.\ii[Key step.] Treat the range of $T$ as the intermediate space.\end{motive}
\begin{prob}\begin{proof}\begin{steps}\ii The map $T:V\to\im T$ is onto by definition and injective by Proposition 32, hence invertible by Proposition 29.\ii The restriction of $S$ is onto because every $v=STu$ has a preimage $Tu\in\im T$. If $w=Tu$ lies in its kernel, then $STu=0$, so $u=0$ and $w=0$. Thus it too is invertible.\ii The same two-composition calculation as Proposition 32, with these domains, yields the displayed formula.\end{steps}\end{proof}\end{prob}
\begin{example*}[Notes Example 37]On $V=\mathcal P(\mathbb R)$, let $Sp=p'$ and $(Tp)(x)=\int_0^{x+1}p(t)\,dt$.
\begin{steps}\ii Termwise integration and differentiation give $STp(x)=p(x+1)$; its inverse sends $p(x)$ to $p(x-1)$.
\ii Every $Tp$ vanishes at $x=-1$. Conversely, if $q(-1)=0$, set $p(x)=q'(x-1)$; then $Tp(x)=q(x)-q(-1)=q(x)$. Thus $\im T=U=\{q:q(-1)=0\}$.
\ii The inverse of $S|_U$ sends $p$ to $\int_{-1}^x p(t)\,dt$; the inverse of $T:V\to U$ sends $q$ to $q'(x-1)$. Composing these gives $p(x-1)$ as required.
\ii $TS$ sends every constant to zero, so it is not invertible, even though $ST$ is invertible.\end{steps}\end{example*}
\begin{proposition*}[Notes Proposition 33; LADR 3D.11]For operators $S,T$ on a finite-dimensional space, $ST=I$ iff $TS=I$.\end{proposition*}
\begin{motive}\ii[Want.] Check only one product when verifying a square inverse.\ii[Know.] Proposition 32 makes both factors invertible when their product is invertible.\ii[Key step.] Multiply by the inverse already guaranteed by the dimension argument.\end{motive}
\begin{prob}\begin{proof}\begin{steps}\ii If $ST=I$, Proposition 32 gives invertibility of $T$, so $S=S(TT^{-1})=(ST)T^{-1}=T^{-1}$. Hence $TS=I$.\ii Interchange $S,T$ to prove the converse.\end{steps}\end{proof}\end{prob}
\begin{example*}[Notes Example 38]On $\mathcal P(\mathbb R)$, let $Sp=p'$ and $Tp(x)=\int_0^x p(t)\,dt$.
\begin{steps}\ii Termwise calculus gives $STp=p$, so $ST=I$.
\ii It also gives $TSp=p-p(0)$; in particular $TS(1)=0\ne1$. Thus $TS\ne I$ in this infinite-dimensional example.\end{steps}\end{example*}
\begin{shortcuts}\ii For square finite-dimensional matrices, one inverse product equal to $I$ suffices by Proposition 33.\ii When a product is invertible, its first-applied factor must be injective and its last-applied factor must be onto.\end{shortcuts}\begin{trapbox}\ii Invertibility of $ST$ does not force invertibility of $TS$ in infinite dimension: Example 37 is a counterexample.\ii Inverses reverse order. For $S=\left(\begin{smallmatrix}1&1\\0&1\end{smallmatrix}\right)$, $T=\left(\begin{smallmatrix}2&0\\0&1\end{smallmatrix}\right)$ over $\mathbb R$, $S^{-1}T^{-1}=\left(\begin{smallmatrix}1/2&-1\\0&1\end{smallmatrix}\right)$ while $(ST)^{-1}=\left(\begin{smallmatrix}1/2&-1/2\\0&1\end{smallmatrix}\right)$.\end{trapbox}
\section{Inverse matrices and change of basis}\begin{qbox}How can one operator have different matrices? A change of coordinates changes the input and output descriptions, while the operator itself stays fixed.\end{qbox}
\begin{prob}\begin{definition*}[LADR 3.82]For ordered bases $B,C$ of the same space, set $Q=M_C^B(I)$, whose columns are $[b_j]_C$; then $[v]_C=Q[v]_B$.\end{definition*}\end{prob}
\begin{proposition*}[LADR 3.71, 3.81--3.82, 3.84, 3.86]For a basis $B$ of finite-dimensional $V$ and $C$ of finite-dimensional $W$, the map $T\mapsto M_C^B(T)$ is an isomorphism $L(V,W)\to F^{m,n}$, where $n=\dim V$, $m=\dim W$. Now let $B,C$ both be bases of $V$. The coordinate transition $Q=M_C^B(I)$ is invertible with $Q^{-1}=M_B^C(I)$. For $T\in L(V)$, if $A=M_B^B(T)$ and $D=M_C^C(T)$, then $A=Q^{-1}DQ$. For an invertible $T:V\to W$, with $B$ a basis of $V$ and $C$ a basis of $W$, $M_B^C(T^{-1})=(M_C^B(T))^{-1}$.\end{proposition*}
\begin{motive}\ii[Want.] Keep the bases aligned when composing or reversing maps.\ii[Know.] Week 3 proved matrix correspondence is a linear bijection and that matrices of compositions multiply with the middle basis shared.\ii[Key step.] Write the same coordinate output by first changing coordinates and by first applying the operator.\end{motive}
\begin{prob}\begin{proof}\begin{steps}\ii The correspondence is an isomorphism by its proved linear bijectivity and Proposition 29. The two identity compositions give $M_B^C(I)M_C^B(I)=I$ and the reverse identity, proving the transition inverse.\ii For any $v$, $[Tv]_C=Q[Tv]_B=QA[v]_B$, but also $[Tv]_C=D[v]_C=DQ[v]_B$. Every column $[v]_B$ is possible, so $QA=DQ$. Multiply by $Q^{-1}$ to get the formula.\ii Apply the composition rule to $T^{-1}T=I_V$ and $TT^{-1}=I_W$, with the stated domain/codomain bases. These give both matrix inverse identities.\end{steps}\end{proof}\end{prob}
\begin{proposition*}[LADR 3.68, 3.80: inverse identities]If $\dim V=\dim W<\infty$, $T:V\to W$ and $S:W\to V$ are linear, then $ST=I_V$ iff $TS=I_W$. Also $(A^{-1})^{-1}=A$, and $(AB)^{-1}=B^{-1}A^{-1}$ for invertible square matrices.\end{proposition*}
\begin{motive}\ii[Want.] Extend the one-sided criterion to maps between different equal-dimensional spaces.\ii[Know.] A left inverse makes $T$ injective, and equal finite dimensions then give invertibility.\ii[Key step.] Identify $S$ with the unique inverse and use the two inverse equations.\end{motive}
\begin{prob}\begin{proof}\begin{steps}\ii If $ST=I_V$ and $Tv=0$, then $v=STv=0$. Proposition 31 makes $T$ invertible, and $S=S(TT^{-1})=(ST)T^{-1}=T^{-1}$. Thus $TS=I_W$; interchanging the maps gives the converse.\ii The two equations defining $A^{-1}$ also say that $A$ is its inverse, giving $(A^{-1})^{-1}=A$. Proposition 32 applied to coordinate maps gives the reverse-order product formula.\end{steps}\end{proof}\end{prob}
\begin{example*}[LADR 3.62--3.64: contrasting maps]\begin{steps}\ii On $\mathbb R^3$, $T(x,y,z)=(-y,x,4z)$ has inverse $(a,b,c)\mapsto(b,-a,c/4)$, verified by both substitutions.
\ii On $\mathcal P(\mathbb R)$, multiplication by $x^2$ is injective by coefficient comparison, but is not onto since the constant $1$ is not a multiple of $x^2$.
\ii On infinite sequences, the backward shift $(a_0,a_1,\ldots)\mapsto(a_1,a_2,\ldots)$ is onto (prepend any coordinate), but not injective (the first coordinate is unrestricted). These show why finite dimension matters.\end{steps}\end{example*}
\begin{example*}[LADR 3.83: coordinate transitions]In $\mathbb R^2$, let $B=((4,2),(5,3))$ and let $C$ be the standard basis.
\begin{steps}\ii The transition $Q=M_C^B(I)=\begin{pmatrix}4&5\\2&3\end{pmatrix}$ has determinant $2$, hence the two columns form a basis, with $Q^{-1}=\begin{pmatrix}3/2&-5/2\\-1&2\end{pmatrix}$.
\ii For $T(x,y)=(x,0)$, $D=M_C^C(T)=\begin{pmatrix}1&0\\0&0\end{pmatrix}$. Thus $M_B^B(T)=Q^{-1}DQ=\begin{pmatrix}6&15/2\\-4&-5\end{pmatrix}$.
\ii Directly, $T(4,2)=(4,0)=6(4,2)-4(5,3)$ and $T(5,3)=(5,0)=\frac{15}2(4,2)-5(5,3)$, verifying the columns.\end{steps}\end{example*}
\begin{shortcuts}\ii Label the direction of every coordinate transition: $Q=M_C^B(I)$ sends $B$-coordinates to $C$-coordinates.\ii To compute in a new basis, use $QA=DQ$ first; solving it for $A$ determines the order of the factors.\end{shortcuts}\begin{trapbox}\ii Changing only the domain basis gives $DQ=\left(\begin{smallmatrix}4&5\\0&0\end{smallmatrix}\right)$ in the example; changing both gives $Q^{-1}DQ=\left(\begin{smallmatrix}6&15/2\\-4&-5\end{smallmatrix}\right)$. The two operations are different.\ii A matrix inverse has reversed domain/codomain bases: its columns express preimages in the original domain basis.\end{trapbox}
\section{Products and external direct sums}\begin{qbox}How are independent components kept in a single vector space? A pair $(v,w)\in V\times W$ records one vector from each space, even when the spaces have no common ambient space.\end{qbox}
\begin{prob}\begin{definition*}[Notes Chapter 8; LADR 3.87]The Cartesian product $V_1\times\cdots\times V_m$ consists of ordered lists $(v_1,\ldots,v_m)$ with $v_j\in V_j$. All spaces use the same field. Addition and scalar multiplication act componentwise.\end{definition*}\end{prob}
\begin{proposition*}[Notes Proposition 34; LADR 3.89]The componentwise operations make $V_1\times\cdots\times V_m$ a vector space. Its zero is $(0,\ldots,0)$ and the negative of $(v_1,\ldots,v_m)$ is $(-v_1,\ldots,-v_m)$.\end{proposition*}
\begin{motive}\ii[Want.] Check all vector-space rules on a product.\ii[Know.] Each component is already a vector space.\ii[Key step.] Equality of ordered lists means equality in every component.\end{motive}
\begin{prob}\begin{proof}\begin{steps}\ii Componentwise sums and scalar multiples remain in the product by closure in each $V_j$. Associativity and commutativity of addition hold in each coordinate.\ii The displayed zero and negative satisfy their identities in every coordinate. The scalar identity, scalar associativity, and both distributive laws hold coordinatewise by the corresponding axiom in $V_j$.\end{steps}\end{proof}\end{prob}
\begin{example*}[Notes Examples 39--40]\begin{steps}\ii $\mathbb R^n$ is the product of $n$ copies of $\mathbb R$: the tuples and componentwise operations coincide.
\ii The concatenation map $\mathbb R^r\times\mathbb R^s\to\mathbb R^{r+s}$ sends $((x_1,\ldots,x_r),(y_1,\ldots,y_s))$ to $(x_1,\ldots,x_r,y_1,\ldots,y_s)$. It is linear; splitting after coordinate $r$ is its linear inverse.
\ii Rebracketing $(v,(w,z))\leftrightarrow((v,w),z)\leftrightarrow(v,w,z)$ and swapping $(v,w)\leftrightarrow(w,v)$ are likewise linear bijections, verified componentwise.\end{steps}\end{example*}
\begin{prob}\begin{definition*}[Notes Chapter 8: coordinate subspaces]Write $V_{(j)}$ for the subspace of the product with all coordinates zero except possibly coordinate $j$. The insertion map $v\mapsto(0,\ldots,v,\ldots,0)$ is linear and bijective onto $V_{(j)}$, with inverse reading coordinate $j$; thus $V_j\cong V_{(j)}$.\end{definition*}\end{prob}
\begin{proposition*}[Notes Proposition 35]$V_1\times\cdots\times V_m=V_{(1)}\oplus\cdots\oplus V_{(m)}$. For this reason the product is called an external direct sum; its factors need not start as subspaces of one common space.\end{proposition*}
\begin{motive}\ii[Want.] Write each tuple as uniquely specified coordinate contributions.\ii[Know.] Week 1 defines directness by uniqueness of summands.\ii[Key step.] Read the $j$th coordinate to recover the $j$th summand.\end{motive}
\begin{prob}\begin{proof}\begin{steps}\ii Every tuple equals $(v_1,0,\ldots,0)+\cdots+(0,\ldots,0,v_m)$, so the subspaces span the product.\ii In any such decomposition, coordinate $j$ of the total equals coordinate $j$ of its $V_{(j)}$ summand. Thus each summand is forced; the sum is direct.\end{steps}\end{proof}\end{prob}
\begin{proposition*}[Notes Proposition 36; LADR 3.92]If all $V_j$ are finite-dimensional, then $\dim(V_1\times\cdots\times V_m)=\sum_{j=1}^m\dim V_j$.\end{proposition*}
\begin{motive}\ii[Want.] Count the degrees of freedom in a tuple.\ii[Know.] Proposition 35 decomposes the product into isomorphic copies of its factors.\ii[Key step.] Use the dimension formula for a finite direct sum from Week 2.\end{motive}
\begin{prob}\begin{proof}\begin{steps}\ii Inserting a finite basis of $V_j$ into coordinate $j$ gives a basis of $V_{(j)}$: spanning and independence hold coordinatewise. Thus $V_{(j)}$ is finite-dimensional with dimension $\dim V_j$.\ii Apply the direct-sum dimension formula to Proposition 35 and add these dimensions.\end{steps}\end{proof}\end{prob}
\begin{proposition*}[Notes Proposition 37; LADR 3.93--3.94]For subspaces $U_1,\ldots,U_m\subseteq V$, the addition map $\Gamma:\prod_jU_j\to\sum_jU_j$, $\Gamma(u_1,\ldots,u_m)=\sum_ju_j$, is linear and surjective. It is injective iff the sum is direct. If the $U_j$ are finite-dimensional, $\dim(\sum_jU_j)\le\sum_j\dim U_j$, with equality iff the sum is direct.\end{proposition*}
\begin{motive}\ii[Want.] Compare a tuple of summands with their sum in the ambient space.\ii[Know.] Week 1 characterizes directness by the only representation of zero; Week 3 gives rank--nullity.\ii[Key step.] The kernel records precisely the nonunique decompositions.\end{motive}
\begin{prob}\begin{proof}\begin{steps}\ii Addition commutes with linear combinations, so $\Gamma$ is linear. Every element of the sum is by definition some $\sum u_j$, proving surjectivity.\ii The kernel consists of tuples with $\sum u_j=0$. It is zero iff every such tuple is zero, exactly the direct-sum criterion.\ii In finite dimension, rank--nullity gives $\dim(\sum U_j)=\sum\dim U_j-\dim\ker\Gamma$. The last term is nonnegative and is zero iff $\Gamma$ is injective.\end{steps}\end{proof}\end{prob}
\method{Construct a product basis}\begin{steps}\ii A basis $v_1,\ldots,v_r$ of $V$ and $w_1,\ldots,w_s$ of $W$ give $(v_1,0),\ldots,(v_r,0),(0,w_1),\ldots,(0,w_s)$ in $V\times W$.\ii Any $(v,w)$ expands using the two coordinate expansions, proving spanning. A zero relation separates into a zero relation in each basis, so every coefficient is zero.\ii For $\mathbb R^2\times\mathbb R^3$, the resulting five vectors have one coordinate $1$ and all others zero, grouped into a pair and a triple.\ii For $\mathcal P_2(\mathbb R)\times\mathbb R^2$ (LADR 3.91), the basis is $(1,0),(x,0),(x^2,0),(0,(1,0)),(0,(0,1))$, where the first three zeros denote the zero vector of $\mathbb R^2$.\end{steps}
\begin{shortcuts}\ii A product dimension is a sum, because component bases contribute separate basis vectors.\ii To check whether a sum is direct, test whether its addition map has zero kernel.\end{shortcuts}\begin{trapbox}\ii $\dim(V\times W)$ is not generally $(\dim V)(\dim W)$: $\mathbb R^2\times\mathbb R^3$ has dimension $5$, not $6$.\ii A product and an internal sum need not have equal dimension: for $U=W=\mathbb R$ inside $\mathbb R$, $U\times W$ has dimension $2$ but $U+W$ dimension $1$.\end{trapbox}
\section{Affine sets and quotient operations}\begin{qbox}When do two translated copies of a subspace represent the same set? The lines $(1,0)+\mathbb R(0,1)$ and $(1,4)+\mathbb R(0,1)$ are equal because their base points differ vertically.\end{qbox}
\begin{prob}\begin{definition*}[Notes Definitions 22--23; LADR 3.95, 3.97]For a subspace $U\subseteq V$ and $v\in V$, the translate $v+U=\{v+u:u\in U\}$ is an affine set parallel to $U$. Two affine sets are parallel if both are translates of the same subspace. The word coset also means a translate $v+U$.\end{definition*}\end{prob}
\begin{example*}[Notes Examples 41--42]\begin{steps}\ii The line $\{(x,y):x+y=1\}$ equals $(1,0)+\operatorname{span}((-1,1))$, because its points are $(1-t,t)$ for arbitrary real $t$.
\ii The polynomials with derivative $x^2$ are $\frac13x^3+\mathbb R$, where $\mathbb R$ denotes the subspace of constant polynomials. Differentiation checks inclusion; subtracting $\frac13x^3$ gives derivative zero, hence a constant by coefficient comparison.
\ii For fixed $U=\mathbb R(0,1)$, the cosets in $\mathbb R^2$ are exactly the vertical lines: $(a,b)+U=\{(a,t):t\in\mathbb R\}$.\end{steps}
\begin{center}\begin{tikzpicture}[scale=.65,>=stealth]\draw[->](-1,0)--(5,0)node[right]{$x$};\draw[->](0,-1)--(0,3)node[above]{$y$};\draw[blue,thick](1,-.7)--(1,2.7);\draw[red,thick](3,-.7)--(3,2.7);\fill(1,1)circle(2pt);\fill(1,2)circle(2pt);\node[below]at(1,0){$a$};\node[below]at(3,0){$b$};\end{tikzpicture}\\Points on the same vertical line belong to one coset of $\mathbb R(0,1)$.\end{center}\end{example*}
\begin{proposition*}[Notes Proposition 38]An affine set is parallel to exactly one subspace.\end{proposition*}
\begin{motive}\ii[Want.] Recover its subspace independently of the chosen base point.\ii[Know.] Every element of a translate differs from its base point by an element of its subspace.\ii[Key step.] Compare two points of the same affine set.\end{motive}
\begin{prob}\begin{proof}\begin{steps}\ii Suppose $A=v+U=x+W$. Since $v\in x+W$, write $v=x+w_0$ for some $w_0\in W$.\ii For $u\in U$, the point $v+u$ is in $x+W$, so $v+u=x+w$ for some $w\in W$. Subtract to get $u=w-w_0\in W$. Thus $U\subseteq W$.\ii Interchange the roles of $U,W$ to obtain the reverse inclusion, hence equality.\end{steps}\end{proof}\end{prob}
\begin{proposition*}[Notes Proposition 39; LADR 3.101]For a fixed subspace $U$, the following are equivalent: $v-w\in U$; $v+U=w+U$; $(v+U)\cap(w+U)\ne\varnothing$. Thus two cosets are either identical or disjoint.\end{proposition*}
\begin{motive}\ii[Want.] Decide equality or overlap without comparing every point.\ii[Know.] $U$ is closed under addition, subtraction and scalar multiplication.\ii[Key step.] Translate the equality of two representative points into a difference in $U$.\end{motive}
\begin{prob}\begin{proof}\begin{steps}\ii If $v-w\in U$, then every $v+u=w+(v-w+u)\in w+U$. Reversing $v,w$ gives equality.\ii If the cosets are equal, they intersect because $v=v+0$ belongs to $v+U$.\ii If a point equals $v+u_1=w+u_2$, then $v-w=u_2-u_1\in U$. This proves all three implications.\end{steps}\end{proof}\end{prob}
\begin{prob}\begin{definition*}[Notes Definition 24; LADR 3.99, 3.102]The quotient $V/U$ is the set of all cosets $v+U$. Define $(v+U)+(w+U)=(v+w)+U$ and $a(v+U)=(av)+U$. These operations must be independent of the representatives $v,w$, as proved next.\end{definition*}\end{prob}
\begin{proposition*}[Notes Definition 24: well-defined operations; LADR 3.103]The quotient operations are well-defined and make $V/U$ a vector space. Its zero vector is the coset $U=0+U$, and the negative of $v+U$ is $-v+U$.\end{proposition*}
\begin{motive}\ii[Want.] Ensure the operations act on cosets rather than on arbitrary choices.\ii[Know.] Proposition 39 identifies equality of cosets with differences in $U$.\ii[Key step.] Show replacing representatives changes the computed result only by a vector in $U$.\end{motive}
\begin{prob}\begin{proof}\begin{steps}\ii If $v_1+U=v_2+U$ and $w_1+U=w_2+U$, then $v_1-v_2,w_1-w_2\in U$. Their sum is in $U$, so $(v_1+w_1)+U=(v_2+w_2)+U$.\ii Also $a(v_1-v_2)\in U$, so $av_1+U=av_2+U$. Both operations are well-defined.\ii Associativity and commutativity of addition, both distributive laws, scalar associativity, and scalar identity follow by performing the corresponding operations on representatives in $V$.\ii $(v+U)+(0+U)=v+U$ and $(v+U)+(-v+U)=0+U$, verifying the asserted zero and negative.\end{steps}\end{proof}\end{prob}
\method{Compute with cosets}\begin{steps}\ii For $U=\mathbb R(0,1)$, $(1,2)+U=(1,4)+U$ because the representative difference $(0,-2)$ lies in $U$.\ii Adding $(1,2)+U$ and $(3,-1)+U$ gives $(4,1)+U=(4,0)+U$. Multiplication by $2$ gives $2((1,2)+U)=(2,4)+U=(2,0)+U$.\ii Distinct first coordinates give distinct cosets by Proposition 39. The zero coset is the entire vertical axis $U$, not merely the point $(0,0)$.\end{steps}
\begin{shortcuts}\ii To compare cosets, subtract their representatives and test membership in $U$.\ii To add or scale cosets, choose the simplest representative; well-definedness guarantees the result is independent of that choice.\end{shortcuts}\begin{trapbox}\ii An affine set need not be a subspace: $x+y=1$ excludes zero.\ii Parallel affine sets need not intersect: the vertical lines $x=1$ and $x=3$ are disjoint.\ii A coset is a set of vectors, not a chosen vector: $(1,2)+U=(1,4)+U$ above although $(1,2)\ne(1,4)$.\end{trapbox}
\section{Problem set}\begin{qbox}Test invertibility, construct bases, and compare cosets using the results above.\end{qbox}
\begin{prob}\textbf{HW4 P1.} For $A=\left(\begin{smallmatrix}a&b\\c&d\end{smallmatrix}\right)$, (a) verify $A\left(\begin{smallmatrix}d&-b\\-c&a\end{smallmatrix}\right)=(ad-bc)I$. (b) Define $\det A=ad-bc$; prove nonzero determinant implies invertibility and derive the inverse. (c) Prove zero determinant implies noninvertibility without general determinant theory.\\\emph{Hint:} Use the two-by-two proposition and its explicit kernel vector.\end{prob}
\begin{prob}\textbf{HW4 P2.} Decide invertibility and find inverses when they exist: $\left(\begin{smallmatrix}2&1\\5&3\end{smallmatrix}\right)$ over $F$, $\left(\begin{smallmatrix}1&i\\i&-1\end{smallmatrix}\right)$ over $\mathbb C$, and $\left(\begin{smallmatrix}1&2\\3&1\end{smallmatrix}\right)$ over $\mathbb F_5$.\\\emph{Hint:} Compute $ad-bc$ in the stated field.\end{prob}
\begin{prob}\textbf{HW4 P3 (LADR 3D.3).} For $T\in L(V)$ with $V$ finite-dimensional, prove equivalence: (a) $T$ is an isomorphism; (b) for every basis $v_1,\ldots,v_n$ of $V$, the list $Tv_1,\ldots,Tv_n$ is a basis of $V$.\\\emph{Hint:} Use injectivity for image independence and surjectivity for image spanning.\end{prob}
\begin{prob}\textbf{HW4 P4.} For matrices $B,C$ with $BC$ defined, (a) prove $\operatorname{rank}(BC)\le\min\{\operatorname{rank}B,\operatorname{rank}C\}$. (b) Give a strict example with both factors of positive rank.\\\emph{Hint:} Use $\im(BC)\subseteq\im B$ and restrict $B$ to $\im C$.\end{prob}
\begin{prob}\textbf{HW4 P5.} For finite-dimensional $V,W$, prove $\dim(V\times W)=\dim V+\dim W$ by constructing a product basis from bases of the factors.\\\emph{Hint:} Use the two coordinate embeddings from Proposition 35.\end{prob}
\begin{prob}\textbf{HW4 P6.} For $U\subseteq V$ a subspace, (a) prove $v+U$ is a subspace iff $v\in U$, in which case $v+U=U$. (b) Prove $v+U=v'+U$ iff $v-v'\in U$.\\\emph{Hint:} For (a) test whether zero belongs; for (b) apply Proposition 39.\end{prob}
\begin{prob}\textbf{HW4 P7 (optional; source wording).} Suppose $U$ is a subspace of finite-dimensional $V$, and $S:U\to V$ is injective. Prove $S$ extends to an invertible operator $T:V\to V$. The source omits the necessary assumption that $S$ is linear; as written the assertion is false. Give a counterexample, then prove the intended statement with linearity added.\\\emph{Hint:} Try $S(x)=x^3$ on $\mathbb R$; for the linear case extend a basis and its independent image list.\end{prob}
\begin{prob}\textbf{Exam practice (practice-midterm1.pdf, Question 3).} Define $T:\mathbb R^3\to\mathcal P_2(\mathbb R)$ by $T(a,b,c)=(a+b)+(b+c)x+(a+c)x^2$. Use domain basis $B=((1,0,0),(1,1,0),(1,1,1))$ and codomain basis $C=(1,x,x^2)$. (a) Find $M_C^B(T)$. (b) Prove $T$ is an isomorphism.\\\emph{Hint:} Compute the three basis images, then test the kernel and apply Proposition 31.\end{prob}\end{document}
